% -------------------------------------------------------
% ASSIGNED TO: \textcolor{green}{Person 5}
% -------------------------------------------------------
% INSTRUCTION: Paste existing Sections 4.5 through 4.5.3 here.
% ADD the new SOTA baseline comparison subsection (marked in red below).
% Convert first-person to third person throughout.
% -------------------------------------------------------

% [Paste Section 4.5 Evaluation Strategy here]
% [Paste Section 4.5.1 Multi-Criteria Ranking via PROMETHEE II here]
% [Paste Section 4.5.2 Statistical Significance (Wilcoxon Signed-Rank Test) here]
% [Paste Section 4.5.3 Automated Biological Verification (GSEApy) here]

The evaluation strategy employs a Multi-Criteria Decision Analysis (MCDA) framework to objectively rank the extracted rules, moving 
beyond a single-metric (accuracy-based) approach \cite{Matejova2025}.

\subsubsection{Multi-Criteria Ranking via PROMETHEE II}
To determine the top 10 rules from the thousands generated by the genetic algorithm, the PROMETHEE II method will be implemented using the 
pymcdm library \cite{Yedjour2024}. Each rule is treated as an alternative and ranked based on three conflicting criteria:
\\
Fidelity: Measured using the F1-score of the rule’s predictions compared to the black-box model’s predictions. This ensures that the rule 
acts as a faithful surrogate \cite{Guidotti2021}.
\\
Comprehensibility: Measured by the Length of the Rule (i.e.,number of conditions). Shorter rules are prioritised to avoid rule explosion 
\cite{Yedjour2024}.
\\
Stability: Calculated by the variance in rule performance across different cancer types in the GDSC dataset \cite{Ye2024}.
\\
\textcolor{red}{For any comparison between two rule sets (for example, baseline versus SAM-regularised,
or evolutionary versus decision-tree extraction), all rules entering that comparison are ranked within a
single combined PROMETHEE II run. This ensures each rule's net outranking flow is computed against a
common reference set, so that the resulting $\phi$ values are directly comparable before the
Mann--Whitney U test is applied. Ranking each rule set in a separate run would render the flows
incommensurable, as PROMETHEE net flows are defined only relative to the alternatives present in a given
run.}
\subsubsection{Statistical Significance \textcolor{red}{(Mann--Whitney U Test)}}
\textcolor{red}{To confirm that differences in rule quality between the baseline and SAM-regularised
models are statistically significant rather than coincidental, a Mann--Whitney U test (equivalently,
the Wilcoxon rank-sum test) is applied to the net outranking flows of the top 10 baseline rules and the
top 10 SAM rules, drawn from the common PROMETHEE II run described above. The Mann--Whitney U test is appropriate because the baseline and
SAM-regularised rule sets comprise distinct rules with no natural one-to-one correspondence; they
therefore constitute two independent samples rather than matched pairs, for which a paired test such
as the Wilcoxon signed-rank test would be inappropriate. A non-parametric test is required because the
distribution of net outranking flows is bounded, derived from a small sample, and cannot be assumed to
follow a normal distribution \cite{Derrac2011}. A $p$-value threshold of 0.05 is used to determine statistical
significance; a result below this threshold is interpreted as evidence that the difference in rule
quality between the baseline and SAM conditions is unlikely to be due to random chance.}

\subsubsection{Automated Biological Verification (GSEApy)}
The genes identified in the top-ranked rules from both the baseline and SAM-regularised models are submitted to over-representation 
analysis using the GSEApy library \cite{Fang2023}, which provides a programmatic interface to the Enrichr and KEGG pathway databases. 
For each rule set, GSEApy tests whether the genes appearing in the extracted rules are statistically enriched in known biological pathways 
relative to a background gene set. The pathway enrichment p-value is used to determine whether the rules align with established oncogenic 
mechanisms, such as apoptosis or DNA repair pathways, or whether they represent statistically arbitrary associations with no biological 
grounding \cite{Subramanian2005}. A significantly higher enrichment score in the SAM-regularised rule set relative to the baseline 
would constitute evidence that improved generalisation produces not only statistically better rules but biologically more meaningful 
ones, directly addressing sub-research question 4 from Section 2.


% -------------------------------------------------------
% NEW CONTENT REQUIRED — SOTA BASELINE COMPARISON
% (Supervisor feedback: compare our method to SOTA approaches)
% -------------------------------------------------------

\subsubsection{SOTA Baseline Comparison: Decision Tree Rule Extraction}

\textcolor{red}{
% ADD: Decision tree as a SOTA rule extraction baseline.
% This is the cleanest way to satisfy the supervisor's request to compare to SOTA
% without adding entirely new algorithm implementations.
% Suggested text:
To contextualise the quality of rules produced by the GA and NSGA-II pipelines, a CART
decision tree \cite{Breiman1984} is included as a deterministic rule extraction baseline,
representing the most widely adopted alternative to evolutionary rule extraction in the XAI
literature \cite{Arrieta2020}. The decision tree is fitted to the predictions of both the baseline
and SAM-regularised DeepCCDS models using scikit-learn, with maximum depth constrained to
prevent rule explosion. Rules are extracted from the fitted tree and evaluated under the identical
PROMETHEE II framework described in Section 4.5.1, using the same fidelity, comprehensibility,
and stability criteria. This enables a direct comparison between deterministic surrogate-based
rule extraction and the proposed evolutionary approach across both model conditions. \textcolor{red}{The Mann--Whitney U test described in Section~\ref{subsec:statsig} is applied to the
PROMETHEE outranking flow distributions of the decision tree rules and the GA/NSGA-II rules to
determine whether the improvement from evolutionary extraction is statistically significant, and
whether this advantage is maintained after SAM regularisation.}
% NOTE: Breiman et al. (1984) "Classification and Regression Trees" is the standard CART citation.
% Add to references.bib: Breiman, L., Friedman, J., Stone, C.J., Olshen, R.A. (1984).
% Classification and Regression Trees. CRC Press.
}